\section{アプローチ09（体の向き取得および距離・列線に基づくグループ判定手法）}

\subsection{手法の概要}
アプローチ09では、アプローチ08までのバウンディングボックス最適化処理によって最終確定した人物領域に対し、身体の向き（姿勢ベクトル）の抽出と、待機列の幾何的構造（列線）に基づくグループ判定アルゴリズムを導入した。

検出された各人物に対して MoveNet Lightning を再適用し、二次元のキーポイント座標から映像平面上の身体の向きベクトルを算出する。これに対応するため、型定義（\texttt{Person} インターフェース）を拡張し、バウンディングボックス情報、骨格座標、および体の向きベクトルを保持する設計とした。Worker Manager においても、人物検出用と1人対象のポーズ検出用（\texttt{processCandidate}）の処理キューを分離・対応させた。

グループ判定においては、個々の人物のバウンディングボックス配置と身体の向きベクトルから待機列全体を表現する1本の近似曲線（列線）を算出するアルゴリズム（\texttt{estimateQueueLine}）を構築した。この列線に対して物理的に横並びに位置する人物同士をグループとして識別するロジック（\texttt{groupByDistance}）を実装し、従来の空間距離のみに依存していた簡易的なグループ化ロジック（\texttt{convertToGroups}）を完全に刷新・代替した。

\subsection{結果およびアプローチ08との比較・課題}

\subsubsection{アプローチ08との比較}
バウンディングボックスの位置に加えて身体の向きと列線の構造を考慮したことで、大きく蛇行する列や非常に長く横方向に伸びる待機列であっても、列の進行方向に対する正確な横並びグループを判定することが可能となった。これにより、列の折り返し部分で前後に隣接する無関係な人物同士を誤って同一グループに統合する問題が解消された。

\subsubsection{判明した課題}
\begin{enumerate}
    \item \textbf{列前後で会話・交差するグループに対する判定の不全}：
    待機列内においては、必ずしも全員が進行方向を向いて綺麗に横並びで並ぶわけではなく、会話やスマートフォンの操作等によって前後の列に跨って向き合うグループが存在する。このような前後の動きや対話関係を持つ連れを検知できない課題が判明した。
    \item \textbf{複数パイプライン（一次検出・ポーズ検出）による計算オーバーヘッド}：
    バウンディングボックスの重複排除および向き取得のために複数回ポーズ検出 Worker や前処理を経由するため、パイプライン全体として依然として計算ステップ数が多く、処理効率のさらなる改善が求められた。
\end{enumerate}
